\section{Introduction}
\label{ch:intro}

The Foundational Reasoning and Feedback Protocol (FRFP) is a formal framework for structuring
reasoning and evaluation in systems that combine machine computation with human judgment. It
is based on a fundamental distinction between \emph{explicit} processes, which operate on
representations and can be executed mechanically, and \emph{tacit} processes, which determine
correctness, validity, and acceptance.

Earlier presentations of FRFP introduced this structure as a conceptual and theoretical framework
for long-horizon human--AI collaboration. This paper provides an updated and consolidated formal
specification of the framework. The presentation reflects refinements arising from complete
formalization, including simplification of the primitive basis, clarification of key definitions,
and strengthening of structural constraints.

The purpose of this document is to present the current form of FRFP as a self-contained system.
The focus is on precise definitions and derivable properties rather than on formal verification or proof
mechanics. Where earlier versions differ, the refined formulation given here is taken as authoritative.

The remainder of the paper develops the framework starting from a minimal categorical kernel
that captures the separation between explicit and tacit domains.

\subsection*{Validation Boundary: Lean and Human Roles}

This theory paper follows the same validation boundary used in the Lean companion.
Lean is the correctness oracle for formal derivations: once a theorem is accepted by
Lean, no human re-check of proof steps is required. Human validation serves a
\emph{different} role: it governs commitments.

Concretely:
\begin{itemize}
  \item \textbf{Lean verifies proofs}: derivations are correct relative to stated definitions and axioms.
  \item \textbf{Humans validate commitments}: statement intent, premise legitimacy, and publication scope.
\end{itemize}

Under FRFP, this distinction is structural rather than optional: explicit derivation is
machine-checkable, while normative acceptance remains a tacit-space responsibility.

\subsection*{The FRFP Document Ecosystem}

This paper is part of a broader research program organized into three layers. Understanding
where this document fits will help the reader navigate the full body of work.

\textbf{Theory layer.} The mathematical foundations of FRFP are developed across three documents.
\textit{FRFPPaper Vol.~1} introduced the explicit--tacit separation, the representational boundary,
and the formal protocol constraints. The present document supersedes that account with a
consolidated and corrected specification, incorporating refinements that emerged from complete
formalization. For readers who require a pedagogical introduction to the underlying mathematics,
the \textit{FRFP Companion Tutorial} provides accessible explanations of the key structures
before any formal development. The \textit{Lean Formalization} companion paper records the
machine-checked verification of this theory in Lean~4, including 269 theorems, a basis reduction
from 11 premises to 7, and novel results not present in earlier versions.

\textbf{Engineering layer.} \textit{FRFPPaper Vol.~2} applies the theoretical framework to a
corpus of canonical socio-technical case studies. The \textit{FRFP Technical Specification}
translates the abstract protocol constraints into implementable procedures for LLM-based
reasoning pipelines. A series of \textit{FRFP White Papers} (e.g., \textit{Beyond Prompt
Engineering}) provides engineering-facing interpretations of the protocol and its practical
consequences.

\textbf{Evaluation layer.} \textit{FRFP Vol.~3: Experimental Methods for Evaluating AI Agents} defines a
general controlled-evaluation framework for AI reasoning systems. The \textit{FRFP vs.\ Base
Survey} applies that framework to empirically assess protocol-guided reasoning against baseline
LLM outputs, using normalized one-page analytical summaries as the evaluation artifacts for
blinded human raters.

\medskip
\noindent\textbf{Recommended reading path.} A reader new to FRFP should begin with
\textit{FRFPPaper Vol.~1}, which provides the motivation, conceptual development, and
theoretical arguments underlying the framework. The present document should then be read
as the authoritative formal specification: it refines and consolidates the definitions from
Vol.~1 but does not repeat its motivating discussion. Readers who find the mathematics
unfamiliar should consult the FRFP Companion Tutorial before or alongside this document.
After this document, readers interested in machine-checked proofs should proceed to the
Lean Formalization companion. Readers whose primary interest is deployment or empirical
evaluation should continue to the Engineering and Evaluation layers in sequence.

Practitioners and engineers whose primary interest is implementation rather than formal
theory may proceed directly to \textit{FRFPPaper Vol.~2}, which provides an
engineering-oriented exposition of the framework through detailed socio-technical case
studies, and to the \textit{FRFP Technical Specification} for operational protocol rules.
